% GENERATED FILE. Edit canonical lessons, then run npm run generate:books.
% canonical-source-hash: fnv1a64:4f4998ec1f11ce2f
% canonical-lessons: HI-C194-dhobi, HI-C194-bavarci, HI-C194-draivar, HI-C194-injiniyar, HI-C194-lekhak

\chapter{A washerman, A cook, A driver, An engineer, A writer}
\label{ch:hi-a-washerman-a-cook-a-driver-an-engineer-a-writer}
\hlchaptermodality{194}

\begin{chapteropening}
\textbf{By the end of this chapter:} \emph{I can say a washerman, a cook, a driver, an engineer and a writer.}

Complete the last lesson of chapter 194: I can say a washerman, a cook, a driver, an engineer and a writer.
\end{chapteropening}

\section[\texorpdfstring{\hi{धोबी} (dhobī)}{धोबी (dhobī)}]{\hi{धोबी} (dhobī) — a washerman}

\label{lesson:HI-C194-dhobi}



\begin{quote}
Before the new one: say the Hindi for a tailor, then the Hindi for a barber.
\end{quote}

\subsection*{You'll want to know: \hi{धोबी}}
\textbf{\hi{धोबी}} — \emph{dhobī} — \textquotedblleft{}a washerman\textquotedblright{}.

\begin{cousinweb}
\textbf{\hi{धोबी}} (\emph{dhobī}) is a tadbhava, a word inherited and reshaped by sound change. It comes from \emph{*dhauvati}, \textquotedblleft{}washes\textquotedblright{}, a form CDIAL reconstructs from its descendants (CDIAL 6886). CDIAL notes that the forms with \emph{-b-} began in the east and spread into every language around.
\end{cousinweb}

\begin{grammarlens}[title={What you've built}]
One more word for this chapter.
\end{grammarlens}

\subsection*{Your turn}
\begin{itemize}
\raggedright
  \item \emph{Say it:} \emph{dhobī}
  \item \emph{Say it:} \emph{dhobī}, once more
  \item \emph{Say it:} say \emph{dhobī}
  \item \emph{Recall:} say the Hindi for a tailor, then the Hindi for a barber, then say \emph{dhobī} again
\end{itemize}

\begin{tcolorbox}[breakable,colback=teal!4,colframe=teal!35!black,title={Before you move on}]
What does \hi{धोबी} mean? (\textquotedblleft{}A washerman\textquotedblright{}.) Say it once more.
\end{tcolorbox}
\section[\texorpdfstring{\hi{बावर्ची} (bāvarcī)}{बावर्ची (bāvarcī)}]{\hi{बावर्ची} (bāvarcī) — a cook}

\label{lesson:HI-C194-bavarci}



\begin{quote}
Before the new one: say the Hindi for a barber, then the Hindi for a washerman.
\end{quote}

\subsection*{You'll want to know: \hi{बावर्ची}}
\textbf{\hi{बावर्ची}} — \emph{bāvarcī} — \textquotedblleft{}a cook\textquotedblright{}.

\begin{cousinweb}
\textbf{\hi{बावर्ची}} (\emph{bāvarcī}) is a loan, usually traced to Persian \emph{bāwarchī}, \textquotedblleft{}cook\textquotedblright{}.
\end{cousinweb}

\begin{grammarlens}[title={What you've built}]
One more word for this chapter.
\end{grammarlens}

\subsection*{Your turn}
\begin{itemize}
\raggedright
  \item \emph{Say it:} \emph{bāvarcī}
  \item \emph{Say it:} \emph{bāvarcī}, once more
  \item \emph{Say it:} say \emph{bāvarcī}
  \item \emph{Recall:} say the Hindi for a barber, then the Hindi for a washerman, then say \emph{bāvarcī} again
\end{itemize}

\begin{tcolorbox}[breakable,colback=teal!4,colframe=teal!35!black,title={Before you move on}]
What does \hi{बावर्ची} mean? (\textquotedblleft{}A cook\textquotedblright{}.) Say it once more.
\end{tcolorbox}
\section[\texorpdfstring{\hi{ड्राइवर} (ḍrāivar)}{ड्राइवर (ḍrāivar)}]{\hi{ड्राइवर} (ḍrāivar) — a driver}

\label{lesson:HI-C194-draivar}



\begin{quote}
Before the new one: say the Hindi for a washerman, then the Hindi for a cook.
\end{quote}

\subsection*{You'll want to know: \hi{ड्राइवर}}
\textbf{\hi{ड्राइवर}} — \emph{ḍrāivar} — \textquotedblleft{}a driver\textquotedblright{}.

\begin{cousinweb}
\textbf{\hi{ड्राइवर}} (\emph{ḍrāivar}) is a loan from English \emph{driver}.
\end{cousinweb}

\begin{grammarlens}[title={What you've built}]
One more word for this chapter.
\end{grammarlens}

\subsection*{Your turn}
\begin{itemize}
\raggedright
  \item \emph{Say it:} \emph{ḍrāivar}
  \item \emph{Say it:} \emph{ḍrāivar}, once more
  \item \emph{Say it:} say \emph{ḍrāivar}
  \item \emph{Recall:} say the Hindi for a washerman, then the Hindi for a cook, then say \emph{ḍrāivar} again
\end{itemize}

\begin{tcolorbox}[breakable,colback=teal!4,colframe=teal!35!black,title={Before you move on}]
What does \hi{ड्राइवर} mean? (\textquotedblleft{}A driver\textquotedblright{}.) Say it once more.
\end{tcolorbox}
\section[\texorpdfstring{\hi{इंजीनियर} (injīniyar)}{इंजीनियर (injīniyar)}]{\hi{इंजीनियर} (injīniyar) — an engineer}

\label{lesson:HI-C194-injiniyar}



\begin{quote}
Before the new one: say the Hindi for a cook, then the Hindi for a driver.
\end{quote}

\subsection*{You'll want to know: \hi{इंजीनियर}}
\textbf{\hi{इंजीनियर}} — \emph{injīniyar} — \textquotedblleft{}an engineer\textquotedblright{}.

\begin{cousinweb}
\textbf{\hi{इंजीनियर}} (\emph{injīniyar}) is a loan from English \emph{engineer}.
\end{cousinweb}

\begin{grammarlens}[title={What you've built}]
One more word for this chapter.
\end{grammarlens}

\subsection*{Your turn}
\begin{itemize}
\raggedright
  \item \emph{Say it:} \emph{injīniyar}
  \item \emph{Say it:} \emph{injīniyar}, once more
  \item \emph{Say it:} say \emph{injīniyar}
  \item \emph{Recall:} say the Hindi for a cook, then the Hindi for a driver, then say \emph{injīniyar} again
\end{itemize}

\begin{tcolorbox}[breakable,colback=teal!4,colframe=teal!35!black,title={Before you move on}]
What does \hi{इंजीनियर} mean? (\textquotedblleft{}An engineer\textquotedblright{}.) Say it once more.
\end{tcolorbox}
\section[\texorpdfstring{\hi{लेखक} (lekhak)}{लेखक (lekhak)}]{\hi{लेखक} (lekhak) — a writer}

\label{lesson:HI-C194-lekhak}



\begin{quote}
Before the new one: say the Hindi for a driver, then the Hindi for an engineer.
\end{quote}

\subsection*{You'll want to know: \hi{लेखक}}
\textbf{\hi{लेखक}} — \emph{lekhak} — \textquotedblleft{}a writer\textquotedblright{}.

\begin{cousinweb}
\textbf{\hi{लेखक}} (\emph{lekhak}) is Sanskrit \emph{lekhaka}, \textquotedblleft{}scribe\textquotedblright{}, taken over whole as a learned loan, a tatsama (CDIAL 11102). It is built on \emph{likh-}, \textquotedblleft{}to write\textquotedblright{}, the root of \textbf{\hi{लिखना}}.
\end{cousinweb}

\begin{grammarlens}[title={What you've built}]
One more word for this chapter.
\end{grammarlens}

\subsection*{Your turn}
\begin{itemize}
\raggedright
  \item \emph{Say it:} \emph{lekhak}
  \item \emph{Say it:} \emph{lekhak}, once more
  \item \emph{Say it:} say \emph{lekhak}
  \item \emph{Recall:} say the Hindi for a driver, then the Hindi for an engineer, then say \emph{lekhak} again
\end{itemize}

\begin{tcolorbox}[breakable,colback=teal!4,colframe=teal!35!black,title={Before you move on}]
What does \hi{लेखक} mean? (\textquotedblleft{}A writer\textquotedblright{}.) Say it once more.
\end{tcolorbox}
